\chapter{System Design and UML Diagrams}

This chapter presents the design of Nexis at four complementary levels of
abstraction. Section~4.1 gives the physical view: the service topology of the
deployed system. Section~4.2 gives the functional view through Level~0 and
Level~1 use case diagrams, covering the three human roles and the nine
autonomous agents. Section~4.3 gives the data view through a context-level
data flow diagram and two Level~1 decompositions --- one for the Owner/Admin
control surface and one for the self-healing loop itself. Section~4.4 closes
with a UML sequence diagram that traces a single complete recovery run,
message by message, from fault detection to post-deploy verification. The
self-healing loop is deliberately given the most space, because it is the
central contribution of the project.

\section{System Architecture Overview}

Figure~\ref{fig:architecture} shows the full service topology of Nexis. The
diagram is organised into horizontal layers, read top to bottom: clients, the
control-plane together with the LLM spine, the sidecar services, the data and
state layer, and the external integrations.

\begin{figure}[htbp]
\centering
\includegraphics[width=\textwidth]{../assets/diagrams/architecture.png}
\caption{NEXIS System Architecture ( full services topology across the client,
control-plane/LLM, sidecar, data, and integration layers ).}
\label{fig:architecture}
\end{figure}

\textbf{Clients.} The topmost layer contains every entry point into the
platform: the web console (\code{apps/web}, Next.js~16), inbound GitHub
webhooks, inbound Sentry/Datadog/PagerDuty alert webhooks
\cite{sentry,datadog,pagerduty}, Slack interactivity callbacks
\cite{slackapi}, and plain HTTP clients (CLI/\code{curl}) that address the
deploy-engine and the public API directly. All of these converge on the
control-plane; no client talks to a data store or an external integration
directly.

\textbf{Control-plane and LLM spine.} The second layer is the heart of the
system. The control-plane is a single Go service (\code{:8080}) that hosts the
HTTP API (chi router, with authentication, RBAC, row-level-security binding,
and rate limiting), the always-on Sentinel detector goroutine (EWMA $+3\sigma$
spike detection), and the cron scheduler (daily digest, QA continuous loop,
billing tickers, schema-drift detection). Inside it, the Temporal
RecoveryPipeline workflow \cite{temporal} orchestrates the nine-agent
self-healing loop as a durable workflow: detect $\rightarrow$ diagnose
$\rightarrow$ plan $\rightarrow$ design $\rightarrow$ code $\rightarrow$ test
$\rightarrow$ package/migrate $\rightarrow$ approve (severity-routed)
$\rightarrow$ deploy $\rightarrow$ verify $\rightarrow$ rollback if needed.
Alongside it sits the LLM spine: OpenAI (gpt-4o family) as the production
default and Ollama \cite{ollama} (local \code{gpt-oss:20b} and
\code{qwen2.5-coder:14b}) as a fully supported alternate provider, a pgvector
retrieval store \cite{pgvector} that supplies code-context passages to the
Architect, Backend, and QA agents, and a structured-output stage that
validates every agent completion against a per-agent JSON Schema and retries
on mismatch.

\textbf{Sidecar services.} The third layer holds the independently deployable
modules. Four are Go services and one is Python:

\begin{itemize}[itemsep=2pt]
  \item \textbf{validator} (\code{:8081}) --- the Docker-sandboxed test runner
        \cite{docker} that shadow-executes every candidate patch under
        \code{--network=none}, \code{--read-only}, \code{--cap-drop=ALL} and
        runs the QA-generated Hypothesis property tests \cite{hypothesis};
  \item \textbf{gitops} (\code{:8082}) --- GitHub App pull-request automation
        (blob $\rightarrow$ tree $\rightarrow$ commit $\rightarrow$ ref
        $\rightarrow$ PR), bearer-authenticated from the control-plane;
  \item \textbf{causal-inference} (\code{:8090}) --- the Python/FastAPI
        graph-evidence ranker that scores Pathfinder's candidate root causes
        by evidence-chain specificity, structural proximity, and textual
        overlap (an honest statistical proxy rather than a fitted structural
        causal model \cite{dowhy});
  \item \textbf{deploy-engine} (\code{:8091}) --- clone, stack detection,
        Dockerfile generation, build, and health-checked container run for
        preview deployments;
  \item \textbf{argocd client} --- the Sync/Rollback integration \cite{argocd}
        used by the post-deploy SLO probe.
\end{itemize}

\textbf{Data and state layer.} The fourth layer contains PostgreSQL~17 with
pgvector, made multi-tenant through row-level security policies \cite{pgrls}
across 33 migrations; Neo4j \cite{neo4j} holding the code dependency graph
that Pathfinder walks; Temporal's own persistence for workflow durability;
Redis for rate limiting and caching; and MinIO/S3 as the encrypted patch
store. The key tenant-scoped tables --- \code{users}, \code{org\_members},
\code{projects}, \code{incidents\_raw}, \code{workflow\_runs},
\code{approval\_decisions}, \code{deployments}, \code{lineage\_events},
\code{audit\_log}, \code{feedback\_examples}, \code{role\_recommendations},
and \code{digest\_reports} --- are listed in the diagram and reappear as the
data stores of the DFDs in Section~4.3.

\textbf{External integrations.} The bottom layer lists every third-party
system Nexis touches: the GitHub App, Sentry, Datadog, PagerDuty, Slack,
ArgoCD, Stripe billing \cite{stripe}, and SES/SMTP mail. Observability across
all layers follows the OpenTelemetry pipeline \cite{otel} into Prometheus,
Loki, and Tempo, visualised in Grafana, and infrastructure is declared as
Terraform code targeting AWS ECS behind an ALB.

The architecture directly reflects the polyglot stack described in
Chapter~3: Go for every latency- and concurrency-sensitive service
(control-plane, validator, gitops, deploy-engine), Python for the
statistics-oriented causal-inference sidecar, and TypeScript for the web
console. The container runtime boundary is the Docker Engine API, which
Podman \cite{podman} satisfies identically --- a property that was exercised
in practice, since the evaluation environment for this report ran the
validator sandbox and deploy-engine on Podman.

\section{Use Case Diagram}
\label{sec:usecase-diagrams}

\subsection{Level 0 Use Case Diagram}

Figure~\ref{fig:usecase-l0} shows the system at its coarsest functional
granularity: a single system boundary, three human actors on the left, and
the external systems and data stores on the right.

\begin{figure}[htbp]
\centering
\includegraphics[width=0.85\textwidth]{../assets/diagrams/usecase-level0.png}
\caption{Level 0 Use Case Diagram ( three human actors and the mediated
external systems ).}
\label{fig:usecase-l0}
\end{figure}

The three human actors correspond to the platform's role-based access model.
The \textbf{Owner/Admin} administers the organisation: membership, roles,
integrations, recovery policy, billing, and --- most importantly --- the
approval gate through which every autonomous repair must pass. The
\textbf{Engineer/Member} performs day-to-day engineering work on the
platform: running demo scenarios, browsing incidents, inspecting agent runs,
and participating in the approval flow. The \textbf{Viewer} is a read-only
experience surfaced through RBAC gating in the console; it can observe
dashboards, incidents, and system health but can change nothing.

On the right-hand side sit the external entities: the GitHub App, the alert
sources (Sentry, Datadog, PagerDuty), Slack, ArgoCD, the LLM provider
(OpenAI or Ollama), Stripe billing, and the platform's own stores
(PostgreSQL/Neo4j and the MinIO patch store). The essential property the
diagram expresses is mediation: every line from an actor terminates at the
system boundary, and every line to an external system originates from it. No
human ever calls GitHub, the LLM, or ArgoCD directly --- the nine autonomous
agents operate entirely inside the boundary and act on the external systems
on the humans' behalf.

\subsection{Level 1 Use Case Diagram}

Figure~\ref{fig:usecase-l1} decomposes the single Level~0 bubble into the
concrete use cases of each actor. The diagram uses colour to group use cases
by the minimum role required: blue for Owner/Admin, orange for the shared
approval flow, green for Engineer/Member, purple for Viewer, and
purple-outlined boxes on the right for the autonomous nine-agent fleet.

\begin{figure}[htbp]
\centering
\includegraphics[width=0.9\textwidth]{../assets/diagrams/usecase-level1.png}
\caption{Level 1 Use Case Diagram ( per-actor use case groups and the
autonomous 9-agent fleet ).}
\label{fig:usecase-l1}
\end{figure}

\textbf{Owner/Admin (blue).} The administrative surface covers fifteen use
cases: managing members and roles; accepting or dismissing role
recommendations produced by the heuristic recommender; connecting
integrations; creating and configuring projects; setting the recovery policy
and SLO; toggling the kill switch; issuing invites and API keys; managing
billing and plan; exporting the RLHF training data (the JSONL feedback
export); viewing the audit log; triggering a manual Sentinel run; exporting
evaluation-harness and audit CSVs; revoking sessions; resetting passwords;
and enrolling or verifying MFA.

\textbf{Shared approval flow (orange).} Four use cases are shared between
Owner/Admin and Engineer/Member, because both roles participate in reviewing
autonomous work: approving, rejecting, or modifying a proposed patch;
watching the live recovery timeline over Server-Sent Events; viewing the
patch diff in the Monaco viewer; and deploying or stopping a preview
container from the project's Ops tab.

\textbf{Engineer/Member (green).} The engineering surface covers running a
Live-Demo scenario, browsing the incidents list, inspecting agent runs,
submitting the NASA-TLX workload survey \cite{nasatlx}, browsing evaluation
harness runs, querying the knowledge base, viewing cost and performance
pages, managing the member's own API keys, updating profile and preferences,
managing the member's own sessions and devices, and the account lifecycle
itself (sign up, sign in, MFA verification).

\textbf{Viewer (purple).} The Viewer holds exactly three use cases:
read-only dashboard, read-only incident view, and read-only system health.

\textbf{The 9-agent fleet (purple-outlined, right).} The fleet appears as a
single autonomous actor whose six grouped use cases summarise the recovery
pipeline: detect and diagnose an incident (Sentinel + Pathfinder); plan and
synthesise the repair (Synthesiser + Architect); generate code and tests
(Backend + QA); package deployment and migrations (DevOps in parallel with
Data Engineer); shadow-validate the patch (Validator sandbox); and open the
pull request, redeploy, or roll back on an SLO breach (GitOps deploy and the
rollback probe). Its only connection to a human use case is the approval
gate --- the fleet executes autonomously inside the boundary and surfaces to
a person exactly once per run.

The permission model is \emph{role-additive}: Owner/Admin holds every
Engineer and Viewer permission, and Engineer holds every Viewer permission.
This is why the diagram draws Owner/Admin connected into the orange shared
group as well as the blue group --- a higher role never loses a capability
that a lower role has, it only gains administrative ones. The read-only
Viewer experience is enforced by RBAC gating in the console even though the
underlying database role enumeration is owner/admin/member, a design
decision discussed in Chapter~3.

\section{Data Flow Diagram}

\subsection{DFD Level 0 (Context Diagram)}

Figure~\ref{fig:dfd-l0} is the context diagram: the entire platform
(control-plane plus the nine-agent loop) as a single process, its external
entities, and the three top-level data stores.

\begin{figure}[htbp]
\centering
\includegraphics[width=0.85\textwidth]{../assets/diagrams/dfd-level0.png}
\caption{DFD Level 0 ( context diagram with external entities and top-level
data stores ).}
\label{fig:dfd-l0}
\end{figure}

The data stores are labelled D1--D3. \textbf{D1, the tenant database
(PostgreSQL)}, receives every organisational record --- users, memberships,
projects, incidents, workflow runs, decisions, and audit entries --- under
row-level security. \textbf{D2, the code graph (Neo4j)}, holds the dependency
graph of each connected repository and is read by Pathfinder during
diagnosis. \textbf{D3, the patch store (MinIO)}, receives the encrypted diff
blob produced by each recovery run.

Four external entity flows cross the boundary. The \textbf{user} (any of the
three roles) sends logins, commands, and approval decisions inward and
receives console views and notifications outward. \textbf{GitHub/GitOps}
receives pull-request creation requests (and Dockerfile reads for the
deploy-engine) and returns PR URLs and commit SHAs. The \textbf{alert
sources} (Sentry, Datadog, PagerDuty, Slack) send fault and alert webhooks
inward --- they are producers of incidents, not consumers of data. The
\textbf{LLM provider} (OpenAI or Ollama) receives completion requests and
returns patch and plan JSON. As at Level~0 of the use case model, every one
of these flows is mediated by the platform process; there is no direct path
from a human to an integration.

\subsection{DFD Level 1 --- Owner/Admin}

Figure~\ref{fig:dfd-l1-admin} decomposes the administrative side of the
platform into twelve numbered processes and the eleven tenant-scoped data
stores they touch.

\begin{figure}[htbp]
\centering
\includegraphics[width=\textwidth]{../assets/diagrams/dfd-level1-admin.png}
\caption{DFD Level 1 --- Owner/Admin ( twelve administrative processes and
their tenant-scoped data stores ).}
\label{fig:dfd-l1-admin}
\end{figure}

Table~\ref{tab:admin-processes} summarises each process and the primary data
store it writes to.

\begin{table}[htbp]
\caption{Owner/Admin DFD Level 1 processes and their primary data stores.}
\label{tab:admin-processes}
\centering
\begin{tabular}{@{}lll@{}}
\toprule
\textbf{Process} & \textbf{Responsibility} & \textbf{Primary store written} \\
\midrule
1.0 Auth + RBAC        & login, JWT session issue, role checks      & \code{users}, \code{sessions} \\
2.0 Members \& Roles   & invite, promote, demote, remove            & \code{org\_members} \\
3.0 Role Recommender   & heuristic promote/downgrade suggestions    & \code{role\_recommendations} \\
4.0 Integrations       & connect GitHub/Slack/Sentry/ArgoCD         & \code{integrations} \\
5.0 Project + Policy   & project config, recovery policy, SLO       & \code{projects} \\
6.0 Approval Gate      & approve / reject / modify decisions        & \code{approval\_decisions} \\
7.0 Billing            & Stripe usage ticker, invoices, plan        & billing tables (Stripe) \\
8.0 Audit \& Export    & hash-chained audit writes, CSV export      & \code{audit\_log} \\
9.0 System Health      & parallel dependency probes (read-only)     & --- (no persistent write) \\
10.0 Sentinel Trigger  & admin-initiated detection run              & \code{incidents\_raw} \\
11.0 Sessions / MFA    & session list/revoke, TOTP enrolment        & \code{sessions} \\
12.0 Kill Switch       & suspend autonomous recovery per project    & \code{workflow\_runs} (halt), \code{audit\_log} \\
\bottomrule
\end{tabular}
\end{table}

Walking the diagram briefly: process 1.0 authenticates the admin and binds
the request-scoped RLS context; 2.0 and 3.0 together form the membership
surface, with the role recommender writing its suggestions to
\code{role\_recommendations} and the accept action flowing back through 2.0
into \code{org\_members}; 4.0 and 5.0 configure what the autonomous loop is
allowed to act on; 6.0 records every terminal approval decision; 7.0 and 8.0
handle billing and the append-only, hash-chained audit trail; 9.0 fans out
health probes to Postgres, Redis, Neo4j, MinIO, and Temporal in parallel and
persists nothing; 11.0 manages sessions and MFA; and 12.0 is the safety
brake that suspends autonomous recovery for a project.

The most important flow in the diagram is process \textbf{10.0, the Sentinel
Trigger}. When an admin manually triggers a detection run, the process
writes directly into \code{incidents\_raw} --- the \emph{same} table the
self-healing loop's automatic detectors write to. There is no separate
``manual'' code path, no parallel queue, and no privileged shortcut: an
admin-triggered run is indistinguishable downstream from an automatic one,
and flows through the identical Temporal pipeline, validator sandbox, and
approval gate. This design decision keeps the system's behaviour uniform and
means every guarantee established for the autonomous path (durability,
validation, severity routing, auditability) holds for manual runs for free.

\subsection{DFD Level 1 --- Self-Healing Loop}

Figure~\ref{fig:dfd-l1-selfheal} decomposes the nine-agent loop itself into
twelve numbered processes, coloured by layer: blue for the detection and
decision agents at the head of the pipeline, green for the Layer-1 execution
team, and orange for the validation, approval, deployment, and verification
tail.

\begin{figure}[htbp]
\centering
\includegraphics[width=\textwidth]{../assets/diagrams/dfd-level1-selfheal.png}
\caption{DFD Level 1 --- Self-Healing Loop ( the 9-agent fleet as twelve
numbered processes with their data stores ).}
\label{fig:dfd-l1-selfheal}
\end{figure}

The pipeline reads end-to-end as follows:

\begin{enumerate}[itemsep=2pt]
  \item \textbf{1.0 Sentinel} receives fault signals --- webhook payloads
        from Sentry/Datadog/PagerDuty or its own EWMA~$+3\sigma$ spike rule
        over the incident stream --- and writes the normalised incident into
        \code{incidents\_raw}.
  \item \textbf{2.0 Pathfinder} walks the Neo4j code graph outward from the
        fault symptom and forwards candidate root-cause nodes, with hop
        distance and degree metadata, to the causal-inference sidecar for
        graph-evidence ranking.
  \item \textbf{3.0 Synthesiser} classifies the incident into a scenario
        (regex fast path with an LLM fallback) and computes the ordered
        \code{selected\_agents} list that the workflow subsequently enforces.
  \item \textbf{4.0 Architect} produces the structured plan ---
        \code{plan\_steps}, \code{affected\_files}, \code{risk\_level} ---
        that acts as the contract every downstream agent is checked against.
  \item \textbf{5.0 Backend} generates the minimal unified diff, guided by
        the LLM provider (shown as a dashed structured-JSON completion flow
        into the green processes).
  \item \textbf{6.0 QA} generates property-based tests from the incident and
        the Backend diff.
  \item \textbf{7.0 DevOps} emits the CI/CD and ArgoCD manifests for the
        change.
  \item \textbf{8.0 Data Engineer} proposes forward and reverse migrations
        where schema changes are involved and emits
        OpenLineage-compatible run events \cite{openlineage} into
        \code{lineage\_events}.
  \item \textbf{9.0 Validator} shadow-executes the patch in the Docker
        sandbox, runs the QA-generated test suite, and stores the encrypted
        candidate diff in the MinIO patch store.
  \item \textbf{10.0 Approval Gate} routes the validated patch by severity
        to the Owner/Admin external entity (the approve/reject/modify flow
        at the bottom left of the diagram) and records the terminal decision
        in \code{approval\_decisions} and \code{feedback\_examples}.
  \item \textbf{11.0 GitOps Deploy} opens the real GitHub pull request
        (writing to the \code{deployments} store and out to the
        GitHub/GitOps/deploy-engine external entity), and redeploys through
        the deploy-engine when the incident's source was a preview
        deployment failure.
  \item \textbf{12.0 Rollback Probe} watches \code{incidents\_raw} during a
        bounded post-deploy window for a fresh fatal on the same
        organisation and service.
\end{enumerate}

The diagram's closing annotation states the property that gives the project
its title: \textbf{the loop is closed}. If process 12.0 observes an SLO
breach after a deployment, it does not raise an alert into a separate
escalation channel --- the breach lands in \code{incidents\_raw} as a fresh
fatal, exactly as any external fault would, and re-enters the pipeline at
process 1.0 as a new Sentinel-detected incident (with an ArgoCD rollback to
the prior known-good revision when the project's
\code{RollbackOnSLOBreach} policy is set). Recovery, verification, and
re-detection form a single cycle over one shared data path.

\section{Self-Healing Loop Sequence Diagram}

Figure~\ref{fig:seq-selfheal} is the UML sequence diagram of one complete
recovery run. All nine agents appear as lifelines, together with the fault
source, the Validator sandbox, the Admin, and the GitOps and Deploy
services. Eighteen numbered messages trace the run in temporal order; each
corresponds to a named Temporal activity, and every one of them is recorded
as a durable ActivityEvent and streamed to the incident timeline over
Server-Sent Events --- which is why the live screenshots in Chapter~5 show
this exact message order as a per-agent timeline.

\begin{figure}[htbp]
\centering
\includegraphics[width=\textwidth]{../assets/diagrams/sequence-selfheal.png}
\caption{Self-Healing Loop Sequence Diagram ( one full recovery run, 18
numbered messages across all nine agent lifelines ).}
\label{fig:seq-selfheal}
\end{figure}

The run proceeds as follows:

\begin{enumerate}[itemsep=2pt]
  \item \textbf{Message 1 --- fault arrival.} A webhook or statistical spike
        (fanned in from any incident source) reaches Sentinel. The resulting
        \code{incidents\_raw} row is what the RecoveryPipeline trigger
        polls, so workflow start is decoupled from signal arrival.
  \item \textbf{Message 2 --- \code{Sentinel.Detect}.} The first workflow
        activity re-records the detection inside the durable workflow. It is
        marked as a stub in the diagram because the detection logic itself
        already ran in the always-on detector goroutine; the activity exists
        so that the workflow timeline owns the provenance of the run.
  \item \textbf{Message 3 --- \code{Pathfinder.Diagnose}.} Pathfinder walks
        the Neo4j dependency graph from the symptom and submits the
        candidates to the causal-inference sidecar for ranking.
  \item \textbf{Message 4 --- \code{Synthesiser.Plan}.} The incident is
        classified into a scenario and the agent set for this run is
        selected; non-selected Layer-1 agents will later appear as
        explicitly ``skipped'' in the timeline rather than silently omitted.
  \item \textbf{Message 5 --- LLM completion.} Drawn as a dashed
        asynchronous flow spanning the agent lifelines: every reasoning
        agent obtains its output as a structured JSON completion that is
        validated against that agent's schema, with a retry on mismatch.
  \item \textbf{Message 6 --- \code{Architect.Solution}.} The structured
        plan with \code{plan\_steps} and \code{affected\_files} is produced;
        this becomes the contract for the rest of the run.
  \item \textbf{Message 7 --- \code{Backend.Codegen}.} The minimal unified
        diff is generated. If the diff touches any file outside the
        Architect's \code{affected\_files}, the workflow flags a contract
        violation, which forces HIGH severity at the approval gate.
  \item \textbf{Message 8 --- \code{QA.TestGen}.} Property-based tests are
        generated from the incident and the diff.
  \item \textbf{Messages 9a/9b --- \code{DevOps.Pipeline}
        $\parallel$ \code{DataEngineer.Migrations}.} The only parallel pair
        in the pipeline: CI/CD and ArgoCD manifests are prepared
        concurrently with the forward and reverse migration proposal, since
        neither depends on the other's output.
  \item \textbf{Message 10 --- hand-off to \code{Validator.Validate}.} The
        Backend result, together with the QA suite, is passed to the
        Validator.
  \item \textbf{Message 11 --- sandboxed execution.} The Validator runs the
        patch in the locked-down container
        (\code{--network=none}, \code{--read-only}, \code{--cap-drop=ALL},
        \code{--pids-limit=128}, \code{--security-opt=no-new-privileges}),
        executes the pytest/Hypothesis suite, and reports pass/fail and
        coverage. No unvalidated patch can proceed past this point.
  \item \textbf{Message 12 --- \code{ApprovalGate.Route}.} Severity is
        computed as a function of the scenario, sensitive-path matches, and
        contract violations: HIGH for the schema-drift scenario, any diff
        touching a sensitive path glob (\code{*.sql}, \code{migrations/},
        \code{auth/}, \code{security/}, \code{crypto/}), an unknown
        scenario, or a contract violation; MEDIUM for everything else; LOW
        for tiny UI-only diffs.
  \item \textbf{Message 13 --- notification.} For HIGH and MEDIUM runs the
        gate notifies humans over Slack and email and waits.
  \item \textbf{Message 14 --- human decision.} The Admin's approve, reject,
        or modify action arrives as a Temporal signal. The subsequent
        \textbf{alt block} encodes the automatic paths: LOW severity
        auto-approves immediately, and MEDIUM auto-approves after a 120-second
        countdown unless a human overrides it first. HIGH always requires an
        explicit human decision --- one of the two live runs evidenced in
        Chapter~5 genuinely timed out at this gate, which is the safety
        mechanism working as designed.
  \item \textbf{Message 15 --- \code{GitOpsDeploy}.} On approval, the gitops
        service performs the blob $\rightarrow$ tree $\rightarrow$ commit
        $\rightarrow$ ref $\rightarrow$ PR sequence against GitHub.
  \item \textbf{Message 16 --- conditional \code{DeployEngine.Redeploy}.}
        Executed only when the triggering incident's source was the
        deploy-engine: a merged fix to a broken preview build triggers an
        automatic rebuild and redeploy of the preview container.
  \item \textbf{Message 17 --- post-deploy health window.} The rollback
        probe watches \code{incidents\_raw} for a fresh fatal on the same
        organisation and service within the bounded window.
  \item \textbf{Message 18 --- conditional \code{ArgoCD.Rollback}.} Inside
        the closing \textbf{alt block}, an SLO breach detected during the
        window triggers the ArgoCD Rollback client to restore the prior
        known-good revision --- and, as established in
        Section~4.3.3, the breach itself re-enters the pipeline as a new
        incident.
\end{enumerate}

The footer of the diagram records the observed steady-state wall-clock time
of a full run: approximately 35 seconds to 2 minutes on local Ollama models,
depending on scenario complexity and LLM latency. This range is corroborated
by the real timings visible in the Chapter~5 screenshots, including the
fully successful run that completed in 1\,m\,40\,s from detection to
verified deployment.

Taken together, the four views of this chapter describe one coherent design:
a mediated platform boundary (use case Level~0), a role-additive human
surface beside an autonomous fleet (use case Level~1), a single shared data
path for manual and automatic operation (DFD Level~1, Owner/Admin), and a
detection--repair--verification cycle that feeds its own failures back into
itself (DFD Level~1 self-healing loop and the sequence diagram). Chapter~5
presents the implementation of each element, and Chapter~6 the evidence that
the loop closes in practice.
